\FloatBarrier\section{Introdução}
Para quem utiliza um serviço que registra operações, receber uma página e conseguir concluir uma operação são experiências distintas. Uma interface pode permanecer acessível enquanto a aplicação perdeu a capacidade de gravar. Essa distinção orientou a análise de uma aplicação web de contagem de visitas: a cada consulta, o cliente consegue observar qual instância respondeu e se a visita foi registrada. O serviço expõe, assim, dois efeitos de uma falha que um teste restrito ao código HTTP poderia confundir.

A conteinerização permite iniciar outra instância do programa a partir da mesma imagem. Para recuperar o serviço, porém, também é preciso detectar a ausência da instância, iniciar sua substituta e decidir quando encaminhar tráfego a ela. Há ainda um problema de estado: se cada cópia mantivesse a contagem apenas em sua memória, remover uma réplica eliminaria parte do histórico e consultar outra produziria um total diferente. No projeto, o Kubernetes gerencia a quantidade de instâncias da aplicação, e o PostgreSQL mantém o registro comum às réplicas em um volume que pode ser remontado após a substituição do pod.

O objetivo geral foi avaliar até onde essa divisão de responsabilidades preserva o atendimento diante de mudanças no cluster. Os objetivos específicos foram verificar a distribuição das chamadas entre réplicas, medir a reposição após uma remoção, observar a expansão manual e automática, acompanhar o contador durante a recriação do banco e testar quais clientes conseguem alcançar sua porta. Os ensaios complementares verificaram o efeito de uma imagem indisponível e a possibilidade de iniciar componentes com imagens já presentes no nó.

A arquitetura é examinada a seguir pelo percurso de uma visita. A análise experimental separa a observação feita pelo cliente da observação dos controladores. O escopo é um laboratório Minikube de nó único: ele permite provocar a perda de um pod conservando o restante do ambiente. Esse desenho define tanto as conclusões possíveis quanto os limites de sua transferência a um serviço distribuído entre máquinas.
